% ============================================================
% Préambule — Rapport PFE
% ============================================================

% --- Langue et encodage ---
\usepackage[french]{babel}
\usepackage[utf8]{inputenc}
\usepackage[T1]{fontenc}
\usepackage{csquotes}

% --- Mise en page ---
\usepackage[a4paper,left=3cm,right=2.5cm,top=2.5cm,bottom=2.5cm]{geometry}
\usepackage{setspace}
\onehalfspacing
\usepackage{fancyhdr}
\usepackage{titlesec}
\usepackage{emptypage}

% --- Couleurs (identité CLEDISS : marine + violet) ---
\usepackage{xcolor}
\definecolor{clednavy}{HTML}{14142B}
\definecolor{cledaccent}{HTML}{A855F7}
\definecolor{cledgray}{HTML}{555566}

% --- Graphiques et figures ---
\usepackage{graphicx}
\usepackage{float}
\usepackage{caption}
\usepackage{subcaption}
\usepackage{tikz}
\usetikzlibrary{shapes.geometric,arrows.meta,positioning,fit,backgrounds,calc,shadows}
\tikzset{every path/.append style={shorten >=1mm, shorten <=1mm}}
\usepackage{pgfgantt}

% --- Tableaux ---
\usepackage{booktabs}
\usepackage{tabularx}
\usepackage{longtable}
\usepackage{multirow}
\usepackage{array}

% --- Listes ---
\usepackage{enumitem}
\setlist{itemsep=2pt, topsep=4pt}

% --- Code source ---
\usepackage{listings}
\lstdefinestyle{code}{
  basicstyle=\ttfamily\footnotesize,
  breaklines=true,
  breakatwhitespace=false,
  frame=single,
  rulecolor=\color{cledgray},
  backgroundcolor=\color{black!3},
  numbers=left,
  numberstyle=\tiny\color{cledgray},
  keywordstyle=\color{cledaccent}\bfseries,
  commentstyle=\color{green!40!black}\itshape,
  stringstyle=\color{orange!70!black},
  showstringspaces=false,
  tabsize=2,
  captionpos=b,
}
\lstset{style=code}

% --- Mathématiques (doit précéder cleveref) ---
\usepackage{amsmath,amssymb}

% --- Liens et références (cleveref doit être chargé en dernier) ---
\usepackage[hidelinks,colorlinks=true,linkcolor=clednavy,citecolor=clednavy,urlcolor=cledaccent]{hyperref}
\usepackage[french]{cleveref}

% --- Titres de chapitres/sections ---
\titleformat{\chapter}[display]
  {\normalfont\huge\bfseries\color{clednavy}}
  {\hfill\Large\textcolor{cledaccent}{\chaptertitlename}~\thechapter}
  {6pt}
  {\hfill}
\titlespacing*{\chapter}{0pt}{-10pt}{30pt}

\titleformat{\section}
  {\normalfont\Large\bfseries\color{clednavy}}{\thesection}{1em}{}
\titleformat{\subsection}
  {\normalfont\large\bfseries\color{clednavy}}{\thesubsection}{1em}{}
\titleformat{\subsubsection}
  {\normalfont\normalsize\bfseries\color{cledgray}}{\thesubsubsection}{1em}{}

% --- En-têtes / pieds de page ---
\setlength{\headheight}{14pt}
\pagestyle{fancy}
\fancyhf{}
\renewcommand{\headrulewidth}{0.4pt}
\fancyhead[C]{\small\textit{\leftmark}}
\fancyfoot[C]{\thepage}

\fancypagestyle{plain}{
  \fancyhf{}
  \fancyfoot[C]{\thepage}
  \renewcommand{\headrulewidth}{0pt}
}

% --- Divers ---
\usepackage{pdfpages}

% Petites commandes utilitaires
\newcommand{\todoimg}[2]{%
  \begin{figure}[H]
    \centering
    \fbox{\parbox[c][4cm][c]{0.85\linewidth}{\centering\itshape #1}}
    \caption{#2}
  \end{figure}%
}

% Petit logo centré, sans légende ni numérotation (pour ne pas polluer la
% liste des figures avec des dizaines de logos). Affiche le fichier s'il
% existe, sinon un encadré de substitution indiquant le nom de fichier
% attendu -- le document compile donc toujours, avant comme après l'ajout
% des vrais fichiers dans images/.
\newcommand{\logofig}[1]{%
  \begin{center}
    \IfFileExists{#1}{%
      \includegraphics[height=1.8cm]{#1}%
    }{%
      \fbox{\parbox[c][1.8cm][c]{4cm}{\centering\itshape Logo manquant\\\texttt{\detokenize{#1}}}}%
    }
  \end{center}%
}
